\section*{What the battery shows}

Across all registrations of 2002--2018 the most leading fifth fails the
five-year proof 1.39 points more often than the most lagging fifth. Six
of the thirty pre-stated contrasts survive adjustment for multiple
testing; four of the six run against the prediction.

\emph{The penalty is concentrated where a class's filing volume is
growing.} Where both the class's theme mix and its filing volume changed
fast, the penalty is 3.2 points; where only volume grew, 1.9; in calm
class-years, 0.6; where only the mix turned over, $-0.2$. The pace-of-change
account predicts the most durable advantage where the market grows and the
technology is stable \citep{SuarezLanzolla2007}; the record puts the
penalty there instead. This is what the crowding reading of lead
predicts: when entrants pour into a class, the language they pour into is
the language that fails.

\emph{Platform and network offerings carry the largest penalty.} Filings
describing marketplaces, platforms, social networks or peer-to-peer
services have a penalty of 7.2 points, against 1.2 for all other filings.
Network effects are listed as a source of first-mover advantage
\citep{LiebermanMontgomery1988}; in the record, early entrants to
networked categories fail more often than late ones.

\emph{In consumer packaged goods the penalty falls on established firms,
not new ones.} Among owners with 25 or more earlier filings, leading
filings fail 4.9 points more often; among first-time filers, 0.9 points
less often. The complementary-assets account predicts the reverse
\citep{Teece1986}: incumbents with distribution should hold what small
pioneers open.

\emph{Manufactured goods carry a larger penalty than services} (2.3
against 1.1 points), the reverse of the shakeout account, in which early
manufacturers are the likeliest survivors \citep{KlepperMiller1995,
Klepper1997}.

\emph{Counsel absorbs most of the penalty.} Self-filed leading filings
fail 4.2 points more often than self-filed lagging ones; counsel-represented,
0.6.

\emph{Two contrasts fall just short of the adjustment} (Holm $p = 0.065$):
the penalty is larger in themes whose share swings from year to year (2.1
against 0.9), as the uncertainty account predicts, and larger for
site-based offerings such as restaurants, hotels and clinics (2.7 against
1.4), against the preemption account.

\emph{Pioneers of a new vocabulary fail, and so do its early followers.}
Filings using a new technology's or product's vocabulary in its first two
years in a class fail 12.8 points more often than other filings of the
same class and year; filings two to five years in, 11.8 points more
often; filings six or more years in, 6.0. The first two do not differ
($t = 0.5$). Golder and Tellis's pioneers failed; here the early followers
failed at the same rate.

\emph{Within verticals the sign differs.} Leading filings in beer fail
6.2 points less often than lagging ones, and in live entertainment 5.6
points less often; in coffee they fail 3.2 points more often and in
fitness 2.9. These vertical estimates are not adjusted for multiple
testing.

\emph{The blue-ocean exemplars' own marks do not lead their classes.}
The mean lead fifth of the registrations of [yellow tail]'s owner,
Cirque du Soleil, Southwest Airlines, NetJets, Accor, Starbucks and Home
Depot lies between 2.7 and 3.2, the middle of the distribution; only
Curves's marks lean leading (3.8). These are the firms' marks of
2002--2018, mostly filed years after each created its category.

\emph{Geography changes where filers fail, not the lead penalty.} A theme's
co-location with itself does not change its penalty (1.8 in the most
concentrated third, 2.3 in the most dispersed). Within themes that have a
hub, filers in the hub fail 1.2 points more often than filers outside it
(s.e.\ 0.45), whatever their lead, and their lead penalty is the same.

\emph{Nothing distinguishes} patenting owners, pharmaceuticals and
devices, consumer against industrial goods, themes imported from another
class, themes grown by first-time filers, invention vocabularies, fads
against surges that held, vocabulary new to its class against new to the
corpus, foreign owners, intent-to-use filings, regulated classes, or
boom against bust years. Several of these groups are small (invention
vocabularies 12,612 registrations; general-purpose-technology filings in
their first three years in a class, 1,166).
